\documentclass[12pt,a4paper]{article}
\usepackage[UTF8]{ctex}
\usepackage{amsmath,amssymb}
\usepackage{booktabs}
\usepackage{geometry}
\geometry{margin=2.5cm}

\title{SPARTA GS 反应微观概率公式（基于 Molchanova et al. 2018）}
\author{}
\date{2026-07-13}

\begin{document}
\maketitle

\section{符号表}
\begin{align*}
  F_N      & \quad \text{每个模拟粒子代表的真实分子数} \quad (\texttt{fnum})\\
  S_e      & \quad \text{表面元素面积} \quad (\texttt{area[isurf]})\\
  \theta_{\max} & \quad \text{最大覆盖度} \quad (\texttt{max\_cover})\\
  N_{\max} & \quad \text{最大可吸附模拟粒子数} \quad (\texttt{maxstick})\\
  N_{\text{tot}} & \quad \text{当前表面吸附总粒子数} \quad (\texttt{total\_state[isurf]})\\
  N_{S}    & \quad \text{空闲表面位数} = N_{\max} - N_{\text{tot}}\\
  N_{AS,k} & \quad \text{吸附物种 }k\text{ 的数目} \quad (\texttt{species\_state[isurf][k]})\\
  V_n      & \quad \text{入射粒子法向速度} = |\mathbf{v} \cdot \mathbf{n}|\\
  T_w      & \quad \text{壁面温度}\\
  k_B      & \quad \text{Boltzmann 常数}\\
  m        & \quad \text{入射粒子质量}
\end{align*}

\section{Arrhenius 反应速率常数}
\begin{equation}
  K(T_w) = A \cdot T_w^{\,b} \cdot \exp\!\left(-\frac{E_a}{k_B T_w}\right)
\end{equation}
对于 SIMPLE 类型反应，$K(T_w)$ 退化为常数：$b = 0,\; E_a = 0$。

\section{GS 反应分类与处理方式}

根据反应是否涉及表面物种参与，将 GS 反应分为两类处理：

\begin{center}
\begin{tabular}{c|c|c}
\toprule
\textbf{类别} & \textbf{反应类型} & \textbf{概率公式} \\
\midrule
常数概率型 & D, E, R & $P = k_{\text{react}}$（与 SPARTA 原实现一致）\\
有限速率型 & AA, DA, LH1, LH3, CD, ER & $P(V_n) = \min\!\left(1,\; \dfrac{2K(T_w)}{V_n} \cdot \Gamma\right)$ \\
\midrule
SPARTA 原样 & CI & 见第\ref{sec:ci}节（与 SPARTA 原实现一致）\\
\bottomrule
\end{tabular}
\end{center}

\section{有限速率统一公式（文献 Eq. 35）}
AA、DA、LH1、LH3、CD、ER 六种反应的微观概率统一为：
\begin{equation}
  \boxed{P(V_n) = \min\!\left(1,\; \frac{2K(T_w)}{V_n} \cdot \Gamma\right)}
\end{equation}
当 $V_n < 10^{-10} \; \text{m/s}$ 时，$P(V_n) = 1$。

其中 $\Gamma$ 为表面密度因子，物理意义为\emph{气相粒子需要与之作用的表面物种的数密度}，
由反应类型决定（见下文）。

\section{D/E/R：常数概率型}

\subsection{DISSOCIATION（解离）}
无表面参与，保持 SPARTA 原有常数概率：
\begin{equation}
  P = k_{\text{react}}
\end{equation}

\subsection{EXCHANGE（交换）}
无表面参与，保持 SPARTA 原有常数概率：
\begin{equation}
  P = k_{\text{react}}
\end{equation}

\subsection{RECOMBINATION（复合消失）}
无表面参与，保持 SPARTA 原有常数概率：
\begin{equation}
  P = k_{\text{react}}
\end{equation}

\section{AA/DA/ER：单因子有限速率型}

\subsection{AA（吸附）}
气相粒子与\emph{空闲位点}作用（A + S $\to$ AS）：
\begin{equation}
  \Gamma = \frac{N_S \cdot F_N}{S_e} \qquad (\text{空闲位点数密度})
\end{equation}
\begin{equation}
  P(V_n) = \min\!\left(1,\; \frac{2 K_{\text{ads}}(T_w)}{V_n} \cdot \frac{N_S \cdot F_N}{S_e}\right)
\end{equation}

\subsection{DA（解离吸附）}
气相分子需\emph{两个空闲位点}（AB + 2S $\to$ A(s) + B(s)）：
\begin{equation}
  \Gamma = \left(\frac{N_S \cdot F_N}{S_e}\right)^2 \qquad (\text{空闲位点数密度的平方})
\end{equation}
\begin{equation}
  P(V_n) = \min\!\left(1,\; \frac{2 K_{\text{disads}}(T_w)}{V_n} \cdot \left(\frac{N_S \cdot F_N}{S_e}\right)^2\right)
\end{equation}

\subsection{ER（Eley-Rideal 复合）}
气相粒子与\emph{吸附态共反应物}作用（AS + B $\to$ AB + S）：
\begin{equation}
  \Gamma = \frac{N_{AS} \cdot F_N}{S_e} \qquad (\text{吸附物种的数密度})
\end{equation}
\begin{equation}
  P(V_n) = \min\!\left(1,\; \frac{2 K_{ER}(T_w)}{V_n} \cdot \frac{N_{AS} \cdot F_N}{S_e}\right)
\end{equation}

\section{LH1/LH3/CD：多因子有限速率型}

这三种反应涉及表面共反应物（可能不止一个），$\Gamma$ 为各表面反应物因子的连乘。

\subsection{表面共反应物状态类型：(s) 与 (b)}

在 LH1/LH3/CD 反应中，表面相反应物的状态标签分为两种：

\begin{center}
\begin{tabular}{c|c|c}
\toprule
\textbf{状态标签} & \textbf{类型} & \textbf{$\Gamma_j$ 处理} \\
\midrule
\texttt{(s)} & 表面吸附态 & 实时计算 $\Gamma_j$（参与反应消耗）\\
\texttt{(b)} & 体相 & $\Gamma_j = 1$（常数，不参与实时更新）\\
\bottomrule
\end{tabular}
\end{center}

体相物种的数密度远大于表面吸附态，其消耗对浓度的相对变化可以忽略，因此设为常数 1。\\
\textbf{注意：}体相参与的烧蚀反应（ablation）\textbf{不在本次修改范围内}，后续单独考虑。

\subsection{表面共反应物因子 $\Gamma_j$}
遍历反应式中所有表面相反应物 $j \ge 1$，\textbf{仅对 (s) 态}按催化性/参与性分别计算：

\begin{equation}
  \Gamma_j =
  \begin{cases}
    \displaystyle \binom{N_{\text{tot}}}{\nu_j} \cdot \left(\frac{F_N}{S_e}\right)^{\nu_j}, & \text{催化性（\texttt{part == 0}}）\\[12pt]
    \displaystyle \binom{N_{AS,k}}{\nu_j} \cdot \left(\frac{F_N}{S_e}\right)^{\nu_j}, & \text{参与性（\texttt{part == 1}}）\\[12pt]
    1, & \text{体相 \texttt{(b)}，不参与计算}
  \end{cases}
\end{equation}
\begin{equation}
  \Gamma = \prod_{j \,\ge\, 1,\; \text{state}\neq\text{(b)}} \Gamma_j
\end{equation}

其中 $\binom{N}{k} = \texttt{stoich\_pow}(N, k)$ 为二项式组合数（SPARTA 已实现）。\\
$\frac{F_N}{S_e}$ 与 SPARTA 原实现保持一致，不额外乘以 $W_s$ 权重。

\subsection{LH1（脱附）}
气相粒子触发\emph{表面共反应物}脱附（A(g) + B(s) $\to$ AB(g)）：
\begin{equation}
  P(V_n) = \min\!\left(1,\; \frac{2 K_{\text{LH1}}(T_w)}{V_n} \cdot \Gamma\right)
\end{equation}

\subsection{LH3（表面物种间反应）}
气相粒子触发与表面共反应物的反应，生成新表面物种：
\begin{equation}
  P(V_n) = \min\!\left(1,\; \frac{2 K_{\text{LH3}}(T_w)}{V_n} \cdot \Gamma\right)
\end{equation}

\subsection{CD（体相消耗）}
表面物种被消耗进入体相：
\begin{equation}
  P(V_n) = \min\!\left(1,\; \frac{2 K_{\text{CD}}(T_w)}{V_n} \cdot \Gamma\right)
\end{equation}

\section{CI（碰撞诱导）}\label{sec:ci}

CI 反应与 SPARTA 原实现保持一致，\textbf{不使用}统一有限速率公式：

\begin{itemize}
  \item 基础形式（无 \texttt{energy} 关键词）：
  \begin{equation}
    P = k_{\text{react}} \cdot \Gamma
  \end{equation}

  \item 含 \texttt{energy} 关键词时，额外引入 $m$、$n$ 两个能量因子参数：
  \begin{equation}
    P = k_{\text{react}} \cdot E_i^{\,m} \cdot (\cos\theta)^{\,n} \cdot \Gamma
  \end{equation}
  其中 $E_i = \frac{1}{2}m|\mathbf{v}|^2$（入射粒子动能），$\cos\theta = V_n / |\mathbf{v}|$（入射角余弦），\\
  $m = \texttt{energy\_coeff[0]}$，$n = \texttt{energy\_coeff[1]}$。
\end{itemize}

上述 $\Gamma$ 与 LH1/LH3/CD 中的定义完全一致（表面共反应物因子的连乘）。

\section{$\Gamma$ 汇总表}

\begin{center}
\begin{tabular}{c|c|l}
\toprule
\textbf{反应类型} & \textbf{$\Gamma$} & \textbf{物理含义} \\
\midrule
D & — & 常数概率 $P = k_{\text{react}}$（SPARTA 原样）\\
E & — & 常数概率 $P = k_{\text{react}}$（SPARTA 原样）\\
R & — & 常数概率 $P = k_{\text{react}}$（SPARTA 原样）\\
\midrule
AA & $\displaystyle \frac{N_S \cdot F_N}{S_e}$ & 空闲位点数密度 \\
DA & $\displaystyle \left(\frac{N_S \cdot F_N}{S_e}\right)^2$ & 空闲位点数密度的平方 \\
LH1 & $\displaystyle \prod \Gamma_j$ & 表面共反应物数密度 \\
LH3 & $\displaystyle \prod \Gamma_j$ & 表面共反应物数密度 \\
CD & $\displaystyle \prod \Gamma_j$ & 表面共反应物数密度 \\
ER & $\displaystyle \frac{N_{AS} \cdot F_N}{S_e}$ & 吸附态共反应物数密度 \\
\midrule
CI & $\displaystyle \prod \Gamma_j$ & 表面共反应物数密度（SPARTA 原样，不用统一公式）\\
\bottomrule
\end{tabular}
\end{center}

\section{概率归一化与反应抽选}
\begin{enumerate}
  \item 计算所有可能反应的概率 $P_i$
  \item $\Sigma = \sum_i P_i$
  \item 若 $\Sigma \le 1$：$P_{\mathrm{scatter}} = 1 - \Sigma$，correction = 1
  \item 若 $\Sigma > 1$：$P_{\mathrm{scatter}} = 0$，correction = $1/\Sigma$
  \item 蒙特卡洛抽选：$r \sim \mathcal{U}(0,1)$；若 $r < P_{\mathrm{scatter}}$ 则散射；否则依次累加 $P_i \cdot \mathrm{correction}$ 选中反应
\end{enumerate}

与 SPARTA 原实现逻辑一致。

\section{实现要点}

\begin{enumerate}
  \item \textbf{D/E/R}：不修改，$P = k_{\text{react}}$ 常数概率。
  \item \textbf{AA/DA/ER}：直接计算单因子 $\Gamma$，代入 $P = \min(1, 2K/V_n \cdot \Gamma)$。
  \item \textbf{LH1/LH3/CD}：遍历表面反应物 $j \ge 1$ 连乘 $\Gamma_j$，代入统一公式。
  \item \textbf{CI}：完全保留 SPARTA 原逻辑 —— 基础形式 $P = k_{\text{react}} \cdot \prod\Gamma_j$，
        energy 开启时额外乘以 $E_i^m \cdot (\cos\theta)^n$。
  \item \textbf{$F_N/S_e$ 权重}：与 SPARTA 原实现一致，不额外乘 $W_s$。
  \item \textbf{概率截断与归一化}：保留 SPARTA 原有逻辑。
\end{enumerate}

\end{document}
